% Generated draft: requires mathematical and source review.
\documentclass[11pt]{article}
\usepackage[margin=25mm]{geometry}
\usepackage{fontspec}
\setmainfont{Latin Modern Roman}
\usepackage{amsmath,amssymb,mathtools,unicode-math}
\setmathfont{Latin Modern Math}
\usepackage{enumitem,needspace,xurl,hyperref}
\setlength{\parindent}{0pt}
\setlength{\parskip}{5pt}
\setlength{\emergencystretch}{3em}
\newcommand{\N}{\mathbb{N}}
\newcommand{\Z}{\mathbb{Z}}
\newcommand{\Q}{\mathbb{Q}}
\newcommand{\R}{\mathbb{R}}
\newcommand{\C}{\mathbb{C}}
\newcommand{\F}{\mathbb{F}}
\newcommand{\A}{\mathbb{A}}
\newcommand{\PP}{\mathbb{P}}
\newcommand{\E}{\mathbb{E}}
\newcommand{\eps}{\varepsilon}
\newcommand{\dd}{\,\mathrm{d}}
\begin{document}
\section*{UnsolvedMath ID: 30002982}
\textbf{Source problem code:} catalogue entry
\subsection{Definitions and mathematical statement}
A nodal curve is a curve whose only singularities are ordinary nodes, locally consisting of two smooth branches meeting transversely. The supplied record does not specify the ground field, projectivity, connectedness, irreducibility, stability, genus, number of nodes, or any generality assumptions on the curves under consideration.

For a nodal curve \(C\), the source names the Kemeny--Wahl map \(kw_C\) and the Halic--Wahl map \(hw_C\). Their constructions, domains, codomains, and required hypotheses are not supplied; these symbols therefore remain names for the source's intended maps rather than defined mathematical operators here.

\textbf{Research question (source formulation).} Can Wahl's conjecture for smooth curves and the Arbarello--Bruno--Sernesi theorem be extended to nodal curves using \(kw_C\), \(hw_C\), or both?

An exact conjecture or theorem cannot be extracted from this record alone. It does not state the versions of Wahl's conjecture and the Arbarello--Bruno--Sernesi theorem to be extended, the admissible class of nodal curves, the map properties that should enter the extension, or the precise geometric conclusion. In particular, no equivalence, implication, or universal quantifier over an unspecified class of curves is asserted in this draft. Specifying these data is part of the unresolved formulation.
\subsection{Short English statement}
The question asks for nodal-curve analogues of Wahl's conjecture and the Arbarello--Bruno--Sernesi theorem, using the Kemeny--Wahl map, the Halic--Wahl map, or both. The supplied record leaves the maps and necessary hypotheses unspecified, so it supports a research question but not a precise quantified conjecture. Its partially solved label is supplied dated curation, not an independently verified conclusion.
\subsection{Sources}
\begin{itemize}
\item Supplied local UnsolvedMath registry record, permanent ID \(30002982\), title \emph{Wahl's Conjecture for Nodal Curves}, collection \emph{Oberwolfach Reports Open Problems}, exact source problem code \texttt{catalogue entry}. The record gives proposed year 2015 and does not identify an individual proposer.
\item Supplied report citation: \emph{Mini-Workshop: Singular Curves on \(K3\) Surfaces and Hyperkähler Manifolds} (2016), Oberwolfach Reports, pp.\ 2939--2967, DOI 10.4171/owr/2015/51. \url{https://doi.org/10.4171/owr/2015/51}. The supplied background describes these reports as documenting mathematical workshops and their open-problem discussions. The citation's 2016 year and the record's 2015 proposed year are retained without reconciliation.
\item Supplied literature reference: Halic, \emph{Modular Properties of Nodal Curves}. \url{https://doi.org/10.1007/s00209-010-0830-2}.
\item Supplied literature reference: Halic, \emph{Surjectivity of the Wahl Map}. \url{https://arxiv.org/abs/1312.3492}.
\item Supplied literature reference: \emph{The Moduli of Singular Curves on K3 Surfaces}. \url{https://doi.org/10.1016/j.jsc.2015.05.001}.
\item Literature assessment dated exactly \texttt{2026-08-21}: the supplied curation labels the problem partially solved and states that natural modified Wahl maps and several non-surjectivity and obstruction results exist in the nodal setting, while a full Wahl-conjecture analogue and its Arbarello--Bruno--Sernesi theorem remain open. This is supplied dated curation derived from the supplied OWR material, not an independently verified literature assessment; its time-sensitive status and cited claims require independent verification.
\item The registry labels the statement \texttt{corrected\_verified} and reports that it was checked against the official Oberwolfach report, preserving that report's hypotheses, notation, and target. This is a supplied verification claim; no source consultation or independent verification is asserted here.
\end{itemize}
\end{document}
